% Part: many-valued-logic
% Chapter: three-valued-logics
% Section: lukasiewicz

\documentclass[../../../include/open-logic-section]{subfiles}

\begin{document}

\olfileid{mvl}{thr}{luk}

\olsection{\L ukasiewicz తర్కం}

బహుమూల్య తర్కానికి ప్రచురితమైన తొలి సవివర ప్రతిపాదనల్లో
ఒకదాన్ని పోలిష్ తత్వవేత్త యాన్‌ \L ukasiewicz 1921లో ఇచ్చాడు.
అరిస్టాటిల్ సముద్రయుద్ధ సమస్య అతనికి ప్రేరణ: “రేపు
సముద్రయుద్ధం జరుగుతుంది” అనే వాక్యం \emph{ఈ రోజు}
నిజమూ కాదనీ అబద్ధమూ కాదనీ అనిపిస్తుంది; దాని సత్యమూల్యం
ఇంకా నిర్ణయించబడలేదు. ఇటువంటి “భవిష్యత్తుపై ఆధారపడిన”
వాక్యాలకు మూడవ సత్యమూల్యాన్ని ప్రవేశపెట్టాలని
\L ukasiewicz ప్రతిపాదించాడు.
\begin{quote}
వచ్చే ఏడాది ఒక నిర్దిష్ట సమయంలో, ఉదాహరణకు డిసెంబరు
21న మధ్యాహ్నం, నేను వార్సాలో ఉంటాననే విషయం ప్రస్తుతం
అనుకూలంగానీ ప్రతికూలంగానీ నిర్ణయించబడలేదని వైరుధ్యం
లేకుండా ఊహించవచ్చు. కనుక ఆ సమయంలో నేను వార్సాలో
ఉండడం సాధ్యమే, కాని తప్పనిసరి కాదు. ఈ పరికల్పనపై,
“వచ్చే ఏడాది డిసెంబరు 21న మధ్యాహ్నం నేను వార్సాలో
ఉంటాను” అనే ప్రతిపాదన ప్రస్తుతం నిజమూ కాదు, అబద్ధమూ
కాదు. అది ఇప్పుడు నిజమైతే, భవిష్యత్తులో నేను వార్సాలో
ఉండడం తప్పనిసరి కావాలి; అది పరికల్పనకు విరుద్ధం.
మరోవైపు అది ఇప్పుడు అబద్ధమైతే, భవిష్యత్తులో నేను
వార్సాలో ఉండడం అసాధ్యం కావాలి; అదీ పరికల్పనకు
విరుద్ధం. అందువల్ల ఈ ప్రతిపాదన ప్రస్తుతం నిజమూ
అబద్ధమూ కాదు; “0” లేదా అసత్యం, “1” లేదా సత్యం
రెండింటికీ భిన్నమైన మూడవ విలువ దానికి ఉండాలి.
ఆ విలువను “$\frac{1}{2}$”తో సూచించవచ్చు. అది
“సాధ్యమైనది”ని సూచించి, “నిజమైనది”, “అబద్ధమైనది”
పక్కన మూడవ విలువగా చేరుతుంది.
\end{quote}
\L ukasiewicz మూడవ సత్యమూల్యానికి మనం $\Undef$ను
వాడతాం.\footnote{\L ukasiewicz ఇక్కడ “సాధ్యం” అనే పదాన్ని
నేడు అసాధారణమైన అర్థంలో, అంటే సాధ్యమే కాని తప్పనిసరి
కాదు అనే అర్థంలో వాడాడు.}

$\lnot$, $\land$, $\lor$ సంయోజకాల సత్యమూల్య ప్రమేయాలను
ఈ వ్యాఖ్యానం నుంచి నిర్ణయించడం తేలిక. భవిష్యత్తుపై
ఆధారపడిన వాక్యానికి నిషేధమూ అలాంటి వాక్యమే; అందువల్ల
$\tf{\lnot}(\Undef) = \Undef$. సంయోగంలో ఒక భాగం
అనిర్ణీతమై, మరొకటి నిజమైతే సంయోగమూ అనిర్ణీతమే:
మొదటి భాగం తరువాత ఎలా తేలుతుందో బట్టి సంయోగం
నిజమో అబద్ధమో కావచ్చు. కాబట్టి \[
    \tf{\land}(\True, \Undef) =
\tf{\land}(\Undef, \True) =
\Undef.
\]
కానీ మరొక భాగం అబద్ధమైతే సంయోగం నిజం కాజాలదు;
అందువల్ల \[\tf{\land}(\False, \Undef) =
\tf{\land}(\Undef, \False) = \False.\]
\sourcecorrection{OLTEMVLLUK-001}{ఆంగ్ల మూలంలో ఈ
రెండవ సమీకరణలో ఒకే ఆర్గ్యుమెంట్ల క్రమం రెండుసార్లు
వచ్చింది. కింది సంయోగ పట్టికకు అనుగుణంగా రెండవ
సంభవంలో క్రమాన్ని తిరగవేశాం.}
మిగతా విలువల్లో, రెండు ఆర్గ్యుమెంట్లూ $\True$ లేదా
$\False$ వంటి నిర్ణీత సత్యమూల్యాలైతే, ఫలితం
సాంప్రదాయిక తర్కంలోలాగే ఉంటుంది.

సోపాధిక సంయోజకం విషయంలో కొంత క్లిష్టత ఉంది.
$q$ భవిష్యత్తుపై ఆధారపడిన వాక్యమనుకుందాం.
$p$ అబద్ధమైతే, $q$ తరువాత ఎలా తేలినా $p \lif q$
నిజమే; కాబట్టి $\tf{\lif}(\False, \Undef) = \True$.
$p$ నిజమైతే, $q$ ఎలా తేలినా $q \lif p$ నిజమే;
కాబట్టి $\tf{\lif}(\Undef, \True) = \True$.
$p$ నిజమైతే $p \lif q$ నిజమో అబద్ధమో కావచ్చు;
అందువల్ల $\tf{\lif}(\True, \Undef) = \Undef$.
అలాగే $p$ అబద్ధమైతే $q \lif p$ నిజమో అబద్ధమో
కావచ్చు; అందువల్ల $\tf{\lif}(\Undef, \False) = \Undef$.
$p$, $q$ రెండూ భవిష్యత్తుపై ఆధారపడిన వాక్యాలుగా
ఉన్న సందర్భం మిగులుతుంది. మొదటి ప్రేరణ ప్రకారం
దానికి నిజానికి $\Undef$ ఇవ్వాలి. అయితే అలా చేస్తే
$!A \lif !A$ \emph{సర్వసత్యం కాదు}. $!A \lor \lnot !A$,
$\lnot(!A \land \lnot !A)$లను వదులుకోవడానికి
\L ukasiewiczకు ఇబ్బంది లేదు; కానీ $!A \lif !A$ను
వదులుకోవడానికి వెనుకంజ వేశాడు. అందుకే
$\tf{\lif}(\Undef, \Undef) = \True$గా నిర్దేశించాడు.

\begin{defn}\ollabel{def:lukasiewicz}
మూడు-విలువల \L ukasiewicz తర్కం కింది మాత్రికతో
నిర్వచించబడుతుంది:
\begin{enumerate}
  \item $\lfalse$, $\lnot$, $\land$, $\lor$, $\lif$
  సంయోజకాలు గల ప్రామాణిక ప్రతిజ్ఞావాక్య భాష $\Lang L_0$.
  \item సత్యమూల్యాల సమితి $V = \{\True, \Undef, \False\}$.
  \item $\True$ మాత్రమే నిర్దేశిత విలువ; అంటే
  $V^+ = \{\True\}$.
  \item $\lfalse$కు $\tf{\lfalse} = \False$.
  మిగిలిన సత్యమూల్య ప్రమేయాలను కింది పట్టికలు ఇస్తాయి:
  \begin{center}
    \begin{tabular}{c|c}
      $\tf{\lnot}$ & \\
      \hline
      $\True$ & $\False$ \\
      $\Undef$ & $\Undef$ \\
      $\False$ & $\True$
    \end{tabular}
    \quad
    \begin{tabular}{c|ccc}
      $\tf{\land}[\LogLuk[3]]$ & $\True$ & $\Undef$ & $\False$ \\
      \hline
      $\True$ & $\True$ & $\Undef$ & $\False$ \\
      $\Undef$ & $\Undef$ & $\Undef$ & $\False$\\
      $\False$ & $\False$ & $\False$ & $\False$
    \end{tabular}
    \\[2ex]
    \begin{tabular}{c|ccc}
      $\tf{\lor}[\LogLuk[3]]$ & $\True$ & $\Undef$ & $\False$ \\
      \hline
      $\True$ & $\True$ & $\True$ & $\True$ \\
      $\Undef$ & $\True$ & $\Undef$ & $\Undef$ \\
      $\False$ & $\True$ & $\Undef$ & $\False$
    \end{tabular}
    \quad
    \begin{tabular}{c|ccc}
      $\tf{\lif}[\LogLuk[3]]$ & $\True$ & $\Undef$ & $\False$ \\
      \hline
      $\True$ & $\True$ & $\Undef$ & $\False$ \\
      $\Undef$ & $\True$ & $\True$ & $\Undef$  \\
      $\False$ & $\True$ & $\True$ & $\True$
    \end{tabular}
  \end{center}
\end{enumerate}
\sourcecorrection{OLTEMVLLUK-002}{ఆంగ్ల మూలం ప్రామాణిక
భాషను పేర్కొన్నా, దానిలోని శూన్యస్థానిక అసత్య స్థిరాంకానికి
విలువ ఇవ్వలేదు. పూర్వ ప్రామాణిక మాత్రిక, తరువాతి మూడు-విలువల
గోడెల్ మాత్రికల్లోని నియమానికి అనుగుణంగా అసత్య విలువను
ఇక్కడ సంపాదకీయంగా స్పష్టం చేశాం. నాలుగు పట్టికలు
మాత్రమే ఆ విలువను నిర్బంధించవు.}
\end{defn}

సులభంగా గమనించగలిగినట్లు, $\lnot$, $\land$,
$\lor$ మాత్రమే కలిగిన \tetoken{సూత్రం}{!!{formula}}~$!A$లోని
అన్ని \tetoken{ప్రతిజ్ఞావాక్య చరాలకు}{!!{propositional variable}s}
$\Undef$ కేటాయిస్తే, ఆ సూత్రపు సత్యమూల్యమూ $\Undef$.
అందువల్ల ఉదాహరణకు సాంప్రదాయిక సర్వసత్యాలు
$p \lor \lnot p$, $\lnot(p \land \lnot p)$
అనేవి $\LogLuk[3]$లో సర్వసత్యాలు కావు;
$\pAssign v(p) = \Undef$ అయినప్పుడు
$\pValue{v}(!A) = \Undef$.

సూత్రంలోని ప్రతిజ్ఞావాక్య చరాలకు $\pAssign v(p) =
\True$ లేదా $\False$ కేటాయించే
\tetoken{సత్యమూల్య కేటాయింపుల్లో}{!!{valuation}s},
$\pValue v(!A)$ సాంప్రదాయిక సత్యమూల్యంతో సరిపోతుంది.

\begin{prop}
  $!A$లోని అన్ని $p$లకు $\pAssign v(p) \in
  \{\True, \False\}$ అయితే,
  $\pValue v(!A)[\LogLuk[3]] = \pValue v(!A)[\LogCL]$.
\end{prop}

\begin{prob}\label{mvl:thr:luk:prob:luk-iff}
  $\LogLuk[3]$లో $\pValue v(!A \liff !B) =
  \pValue v((!A \lif !B) \land (!B \lif
  !A))$ అని నిర్వచించామనుకుందాం. అప్పుడు
  $\liff$ సత్య పట్టిక ఏమిటి?
\end{prob}

చాలా సాంప్రదాయిక సర్వసత్యాలు \LogLuk[3]లోనూ
సర్వసత్యాలే; ఉదాహరణకు $\lnot p \lif (p \lif q)$.
సాంప్రదాయిక తర్కంలోలాగే దీనిని సత్య పట్టికతో
తనిఖీ చేయవచ్చు:
\begin{center}
\begin{tabular}{cc|cccccc}
  $p$ & $q$ & $\lnot$ & $p$ & $\lif$ & $\smash{(}p$ & $\lif$ & $q\smash{)}$ \\ \hline
  \True & \True & \False & \True & \True & \True & \True & \True \\
  \True & \Undef & \False & \True & \True & \True & \Undef & \Undef \\
  \True & \False & \False & \True & \True & \True & \False & \False \\
  \Undef & \True & \Undef & \Undef & \True & \Undef & \True & \True \\
  \Undef & \Undef & \Undef & \Undef & \True & \Undef & \True & \Undef \\
  \Undef & \False & \Undef & \Undef & \True & \Undef & \Undef & \False \\
  \False & \True & \True & \False & \True & \False & \True & \True \\
  \False & \Undef & \True & \False & \True & \False & \True & \Undef \\
  \False & \False & \True & \False & \True & \False & \True & \False \\
\end{tabular}
\end{center}

\begin{prob}
  \LogLuk[3]లో కిందివి సర్వసత్యాలని చూపండి:
  \begin{enumerate}
    \item $p \lif (q \lif p)$
    \item\label{mvl:thr:luk:prob:luk-taut-2} $\lnot(p \land q) \liff (\lnot p \lor \lnot q)$
    \item\label{mvl:thr:luk:prob:luk-taut-3} $\lnot(p \lor q) \liff (\lnot p \land \lnot q)$
  \end{enumerate}
  (\olref[mvl][thr][luk]{prob:luk-taut-2},
  \olref[mvl][thr][luk]{prob:luk-taut-3}లలో
  $!A \liff !B$ను $(!A \lif !B) \land (!B \lif !A)$కు
  సంక్షిప్త రూపంగా తీసుకోండి, లేదా
  \olref[mvl][thr][luk]{prob:luk-iff}కు మీ సమాధానాన్ని వాడండి.)
\end{prob}

\begin{prob}
  కింది సాంప్రదాయిక సర్వసత్యాలు \LogLuk[3]లో
  సర్వసత్యాలు కావని చూపండి:
  \begin{enumerate}
    \item $(\lnot p \land p) \lif q$
    \item $((p \lif q) \lif p) \lif p$
    \item $(p \lif (p \lif q)) \lif (p \lif q)$
  \end{enumerate}
  \sourcecorrection{OLTEMVLLUK-003}{మొదటి సూత్రం
  చివర ఉన్న జతలేని ముగింపు కుండలీకరణను మాత్రమే
  తొలగించాం; సంయోజకాలు, ఉపసూత్రాలు మారలేదు.}
\end{prob}

అన్ని సాంప్రదాయిక సర్వసత్యాలూ $\LogLuk[3]$లో
సర్వసత్యాలు కాకపోయినా, ప్రతి
\tetoken{సత్యమూల్య కేటాయింపులోనూ}{!!{valuation}}
అవి కనీసం $\True$ లేదా $\Undef$ విలువలను
పొందుతాయనుకోవచ్చు. అలా కాదు. ప్రతిదృష్టాంతం:
\[
  \lnot(p \lif \lnot p) \lor \lnot(\lnot p \lif p)
\]
$p$కు $\Undef$ అయితే దీని విలువ $\False$.

\begin{prob}
  \L ukasiewicz తర్కంలో కింది సంబంధాల్లో ఏవి
  నిలుస్తాయి? ప్రతిదానికీ సత్య పట్టిక ఇవ్వండి.
  \begin{enumerate}
    \item $p, p \lif q \Entails q$
    \item $\lnot\lnot p \Entails p$
    \item $p \land q \Entails p$
    \item $p \Entails p \land p$
    \item $p \Entails p \lor q$
  \end{enumerate}
\end{prob}

\L ukasiewicz తన మూడు-విలువల వ్యవస్థ ఆధారంగా
సంభావ్యత తర్కాన్ని నిర్మించాలని ఆశించాడు.
దానికోసం “$!A$ సాధ్యం” అనే ఏకస్థానిక
సంయోజకం $\Diamond !A$ను, దానికి అనుగుణంగా
“$!A$ తప్పనిసరి” అనే $\Box !A$ను
ప్రవేశపెట్టాడు:
\begin{center}
  \begin{tabular}{c|c}
    $\tf{\Diamond}$ & \\
    \hline
    $\True$ & $\True$ \\
    $\Undef$ & $\True$ \\
    $\False$ & $\False$
  \end{tabular}
  \quad
  \begin{tabular}{c|c}
    $\tf{\Box}$ & \\
    \hline
    $\True$ & $\True$ \\
    $\Undef$ & $\False$ \\
    $\False$ & $\False$
  \end{tabular}
\end{center}
మరో మాటలో చెప్పాలంటే, $p$ ఇప్పటికే అబద్ధమని
నిర్ణయించబడకపోతేనే అది సాధ్యం; ఇప్పటికే నిజమని
నిర్ణయించబడితేనే $p$ తప్పనిసరి.

\begin{prob}
  $\Box$, $\Diamond$ల సత్య పట్టికలను చేర్చి
విస్తరించిన~\LogLuk[3]లో
  $\Box p \liff \lnot\Diamond \lnot p$,
  $\Diamond p \liff
  \lnot \Box \lnot p$ సర్వసత్యాలని చూపండి.
\end{prob}

కానీ ప్రతిపాదిత ఈ మోడల్ తర్కపు లోపాలు త్వరలోనే
స్పష్టమయ్యాయి. విషయాలు ఎలా తేలినా
$p \land \lnot p$ నిజం కాజాలదు. కనుక అది
ఇప్పుడు అనిర్ణీతంగా ఉండినా, అసాధ్యమైనదిగా
లెక్కించాలి; అంటే $\lnot \Diamond(p
\land \lnot p)$ సర్వసత్యం కావాలి. అయితే
$\pAssign v(p) = \Undef$ అయితే
$\pValue v(\lnot \Diamond(p \land \lnot p)) =
\False$. సంభావ్యత, అనివార్యత వంటి మోడల్
భేదాలను వ్యక్తీకరించడానికి రెండు సత్యమూల్యాలు
సరిపోవని \L ukasiewicz చెప్పింది సరైనప్పటికీ,
మూడవ సత్యమూల్యాన్ని చేర్చడమూ సరిపోదు.
\sourcecorrection{OLTEMVLLUK-004}{ఆంగ్ల మూలం ఈ
చివరి విలువను అనిర్ణీతంగా ఇచ్చింది. పట్టికల
ప్రకారం చరం అనిర్ణీతమైతే దాని నిషేధం,
వాటి సంయోగం అనిర్ణీతం; సాధ్యత సంచాలకం
ఆ సంయోగానికి సత్యం, బయటి నిషేధం అసత్యం
ఇస్తాయి. అందుకే చివరి విలువను సరిచేశాం.}

\end{document}
